\clearpage
\begin{examplebox}
	\textbf{\textcolor{CExample}{例1（基础）}}
	\textbf{\textcolor{CExample}{题目：}}
	已知函数$f(x)$满足$f(x+4)=f(x)$对任意实数$x$成立，求$f(x)$的一个周期。
	\textbf{\textcolor{CExample}{解题步骤：}}
	\begin{description}[style=nextline, leftmargin=2.8em, labelsep=0.5em, itemsep=0.45em]
		\item[\textbf{第一步（识别条件类型）：}]
			条件为$f(x+a)=f(x)$形式，其中$a=4$。
			\par\textbf{理由：} 这是结论一“加常数型”的标准形式。
		\item[\textbf{第二步（套用周期公式）：}]
			根据结论一，周期$T=|a|$；代入$a=4$得$T=4$。
			\par\textbf{理由：} 直接代入公式$T=|a|$。
		\item[\textbf{第三步（写出结论）：}]
			所以$f(x)$的一个周期是$4$。
			\par\textbf{理由：} 由定义$f(x+4)=f(x)$成立，故$4$是周期。
	\end{description}
	\textbf{\textcolor{CExample}{关键结论：}}
	\textbf{周期$T=4$}

	\medskip
	\textbf{\textcolor{CExample}{例2（稍变形）}}
	\textbf{\textcolor{CExample}{题目：}}
	已知函数$f(x)$满足$f(x+1)=f(x-2)$（$a\neq b$），求$f(x)$的周期。
	\textbf{\textcolor{CExample}{解题步骤：}}
	\begin{description}[style=nextline, leftmargin=2.8em, labelsep=0.5em, itemsep=0.45em]
		\item[\textbf{第一步（改写等式）：}]
			将已知等式视为$f(x+a)=f(x+b)$形式，其中$a=1$，$b=-2$。
			\par\textbf{理由：} 识别为差值型，需正确确定$a$和$b$。
		\item[\textbf{第二步（应用周期公式）：}]
			根据结论四，周期$T=|a-b|$；代入$a=1,b=-2$得$T=3$。
			\par\textbf{理由：} 套用公式$T=|a-b|$。
		\item[\textbf{第三步（验证周期）：}]
			任取$x$，有
			\[
				\begin{aligned}
					f(x+3) &= f((x+2)+1) \\ &= f((x+2)-2) \\ &= f(x),
			\end{aligned}
			\]
			故$3$是周期。
			\par\textbf{理由：} 通过代换验证了周期为$3$。
	\end{description}
	\textbf{\textcolor{CExample}{关键结论：}}
	\textbf{周期$T=3$}

	\medskip
	\clearpage
\textbf{\textcolor{CExample}{例3（综合应用）}}
	\textbf{\textcolor{CExample}{题目：}}
	已知定义在$\mathbb{R}$上的函数$f(x)$满足$f(x+3)=-\frac{1}{f(x)}$（$f(x)\neq0$），且$f(1)=2$，求$f(2026)$的值。
	\textbf{\textcolor{CExample}{解题步骤：}}
	\begin{description}[style=nextline, leftmargin=2.8em, labelsep=0.5em, itemsep=0.45em]
		\item[\textbf{第一步（判断周期类型）：}]
			条件形如$f(x+a)=\pm\frac{1}{f(x)}$，属于倒数型，其中$a=3$，取负号。
			\par\textbf{理由：} 确认属于倒数型周期条件。
		\item[\textbf{第二步（计算周期）：}]
			根据结论三，周期$T=2|a|$；由$a=3$得$T=6$，即$f(x+6)=f(x)$。
			\par\textbf{理由：} 应用周期公式$T=2|a|$。
		\item[\textbf{第三步（转化求值）：}]
			因为$2026=6\times337+4$，所以$f(2026)=f(4)$。令$x=1$，由条件得$f(4)=-\frac{1}{f(1)}$；又$f(1)=2$，故$f(4)=-\frac{1}{2}$。因此$f(2026)=-\frac{1}{2}$。
			\par\textbf{理由：} 利用周期将自变量化小，再用已知条件求函数值。
	\end{description}
	\textbf{\textcolor{CExample}{关键结论：}}
	\textbf{$f(2026)=-\frac{1}{2}$}
\end{examplebox}

\clearpage
\begin{examplebox}
\textbf{例4：第一种分式，先化小自变量再代入}

函数定义在 $\mathbb R$ 上，处处 $f(x)\ne1$，且
\[
 f(x+2)=\frac{f(x)+1}{f(x)-1},\qquad f(0)=2.
\]
求 $f(2026)$。

连续代入两次得周期 $4$。由于 $2026=4\times506+2$，
\[
 f(2026)=f(2)=\frac{f(0)+1}{f(0)-1}
 =\frac{2+1}{2-1}=\boxed{3}.
\]
回验一次：$f(4)=(3+1)/(3-1)=2=f(0)$。

\medskip
\textbf{例5：第二种分式，负数也能正常循环}

函数定义在 $\mathbb R$ 上，处处 $f(x)\ne-1$，且
\[
 f(x+3)=\frac{1-f(x)}{1+f(x)},\qquad f(1)=2.
\]
求 $f(2026)$。

周期为 $6$，$2026=6\times337+4$，因此
\[
 f(2026)=f(4)=\frac{1-f(1)}{1+f(1)}
 =\frac{1-2}{1+2}=\boxed{-\frac13}.
\]
回验：$f(7)=(1+1/3)/(1-1/3)=2=f(1)$。
\end{examplebox}
